\documentclass[11pt]{article}
\usepackage[margin=1in]{geometry}
\usepackage{amsmath,amssymb}
\usepackage{listings}
\usepackage{xcolor}
\usepackage{hyperref}
\hypersetup{colorlinks=true,urlcolor=blue}
\lstset{
  basicstyle=\ttfamily\footnotesize,
  breaklines=true,
  frame=single,
  backgroundcolor=\color{black!3},
  columns=fullflexible,
  keepspaces=true,
}

\title{MorningStar-Lab v1.0 --- Riemann Hypothesis Reconnaissance Certificate}
\author{Entangled Technologies --- The Morning Star Project}
\date{2026-05-25}

\begin{document}
\maketitle

\section*{Scope (honest)}
This certificate records a \emph{numerical reconnaissance} of the
Riemann $\zeta$-function on the critical line $\Re(s)=\tfrac12$,
performed by the MorningStar-Lab \texttt{kernel.py} probe, an
arbitrary-precision \texttt{mpmath} backend running at
\texttt{workdps=50}. It is \textbf{not} a proof of the Riemann
Hypothesis. The Lean~4 axiom-debt-$[\,]$ result that ships in
\texttt{lean-proof/} concerns the M1--M10/M13 BC--CM ($h=1$) spine
and is \emph{independent} of this document.

\section*{Genesis seal}
The append-only ledger \texttt{data/hits.txt} is sealed at construction.
The SHA-256 of the immutable preamble (lines 1--9, ending at
\texttt{--- GENESIS SEAL ---}) is baked into
\texttt{scripts/check-genesis-seal.py} and re-verified before every
probe append.
\begin{lstlisting}
Genesis SHA-256:
  eecbcd9a540aa7a2c90edd23827c73e4d1bb5af641d352f70a5de849b21f875f
Genesis lines 5--9 (frozen):
  437
  1094
  axioms=[] 2026-05-24
  M13_CERT_SHA256=d99b0df4b6e3a5817a90e7f8da44b20261d3b670d31309779a776efbf1bcc668
  --- GENESIS SEAL ---
\end{lstlisting}

\section*{Lean axiom debt}
The frozen M1--M10/M13 spine compiles in \texttt{lean-proof/} with
zero axioms. The current on-disk hash of \texttt{lean-proof/VERIFY.txt}
(regenerated by \texttt{lean-proof/regenerate.sh}) is:
\begin{lstlisting}
sha256(lean-proof/VERIFY.txt) =
  4150d231bb01d1bca0ecd73093002b167be1f4e52439e1dc307fb8f9fdf1bfd8
\end{lstlisting}
The Genesis-pinned \texttt{M13\_CERT\_SHA256=d99b0df4\dots} is the
sealed M13 certificate value at v1.8-BC release time and is
\emph{not} expected to equal the current VERIFY.txt hash --- VERIFY.txt
is regenerated; the Genesis line is immutable.

\section*{Probe contract}
For each probe at $(h, N, \Re(s), \Im(s))$ the kernel writes one
ledger line carrying: timestamp, the four coordinates,
\verb|L_nonvanish|, \verb|RH_ok|, \verb|kms_beta|, a backend tag,
\verb|L_abs|, and a per-line SHA-256 over the canonical body
(including the timestamp, so every probe is uniquely receipted).
The gunsight tolerance is \verb|RH_VANISH_TOL = 1e-12|; a probe with
\verb|L_abs| below this fires \verb|RH_ok=True|.

\section*{First three nontrivial zeros (mpmath, MPMATH\_ZETA)}
Produced by \texttt{python lab.py -c "hunt\_zeros(1, 3)"}. Each zero
location is computed by \texttt{mpmath.zetazero($n$)} at
\texttt{workdps=50}; the kernel then probes at
$(1, 1, 0.5, \mathrm{float}(\Im\,z))$, so each row has its own ledger
SHA. Float64 rounding of the imaginary part means $|L|$ at the
probed point is $\sim 10^{-15}$ rather than $0$, which is honestly
recorded in \verb|L_abs|.
\begin{lstlisting}
n=1  t = 14.13472514173469379
     |L| = 6.6681863422837782815e-16   kms_beta = 2.0
     RH_ok = True   tag = MPMATH_ZETA
     sha = 0bc1fc16c4ab5b09... (this hunt) /
           4d5fad4b6d9343eb... (re-probed by bracket_zero(1, 1e-4))
n=2  t = 21.022039638771554993
     |L| = 1.1609164044328830424e-15   kms_beta = 2.0
     RH_ok = True   tag = MPMATH_ZETA
     sha = 495b8b5f2f0c4ebe...
n=3  t = 25.010857580145688763
     |L| = 8.4988448557528602897e-16   kms_beta = 2.0
     RH_ok = True   tag = MPMATH_ZETA
     sha = 1d20ca18198462fa...
\end{lstlisting}

\section*{Bracket scan around the first zero}
\texttt{python lab.py -c "bracket\_zero(1, 1e-4)"} sweeps
$[t_1-10^{-4}, t_1+10^{-4}]$ in steps of $2\cdot 10^{-5}$. Float
arithmetic happens to land exactly on $t_1 \approx 14.134725141734695$
in the sweep, producing one \verb|RH_ok=True| hit identical to the
center probe:
\begin{lstlisting}
{ "n": 1, "t0": 14.134725141734695,
  "window": 1e-4, "step": 2e-5,
  "zeros_found": [{
    "im": 14.134725141734695,
    "L_abs": "6.6681863422837782815e-16",
    "kms_beta": "2.0",
    "tag": "MPMATH_ZETA",
    "sha": "4d5fad4b6d9343ebf37cd71e91e18be1fd3539e0a3b1e37f2fd5517443b511ed"
  }],
  "zeros_count": 1 }
\end{lstlisting}

\section*{Reproducibility}
\begin{lstlisting}
# Verify the seal:
python scripts/check-genesis-seal.py

# Reproduce the three zero rows above:
python lab.py -c "hunt_zeros(1, 3)"

# Reproduce the bracket scan:
python lab.py -c "bracket_zero(1, 1e-4)"

# Recompute the Lean axiom-debt receipt:
bash lean-proof/regenerate.sh      # requires `lake` on PATH
sha256sum lean-proof/VERIFY.txt
\end{lstlisting}

\section*{What this certificate does \textbf{not} claim}
\begin{itemize}
  \item It does \textbf{not} claim a proof of RH. \verb|RH_ok=True| on
        a probe means $|L(s)| < 10^{-12}$ at that specific
        finite-precision sample, nothing more.
  \item The \texttt{MPMATH\_ZETA} backend is not Lean-verified.
        For $h\!\ge\!2$ the kernel logs \texttt{NEEDS\_SAGE} with
        \verb|reason=h>=2_out_of_scope_for_mpmath_backend|; those
        rows carry a deterministic stub \verb|L_nonvanish=True| and
        must not be read as L-function evaluations.
  \item The names \texttt{hit\_437} and \texttt{hit\_1094} in
        \texttt{lean-proof/TheoremaAureum/AutoLemmas.lean} are
        tautologies (\texttt{True := trivial}); they reference the
        OpenCV cube counts from README Appendix~A and claim nothing
        about number theory.
\end{itemize}

\end{document}
